%! AP Statistics — Unit 8, Lesson 8.5 — Group Activity (Answer Key)
\documentclass[10pt]{article}
\usepackage{apstats-article}
\usepackage{apstats-key}
\newcommand{\TallMath}[1]{$\displaystyle #1\rule[-1.4em]{0pt}{3.2em}$}

\begin{document}

\pageheader{Unit 8, Lesson 8.5}{Group Activity --- Answer Key}
\namepartnerperiod

\vspace{0.08in}

\begin{scenariobox}[Scenario A: Extracurricular Activities at Three Schools]{sky}
A researcher randomly selects 100 students from each of three high schools (Schools A, B, and C)
and asks which extracurricular activity (Sports, Arts, or Neither) each student participates in.
The results are shown below.

\vspace{0.08in}
\renewcommand{\arraystretch}{1.4}
\begin{tabularx}{0.82\linewidth}{l cccc}
    & \textbf{Sports} & \textbf{Arts} & \textbf{Neither} & \textbf{Row Total} \\
    \hline
    \textbf{School A} & 45 & 32 & 23 & 100 \\
    \textbf{School B} & 38 & 40 & 22 & 100 \\
    \textbf{School C} & 52 & 28 & 20 & 100 \\
    \hline
    \textbf{Col Total} & 135 & 100 & 65 & 300 \\
\end{tabularx}
\end{scenariobox}

\begin{enumerate}[label=\textbf{A\arabic*.}]

    \item \textbf{Identify the test type.} Is this a test for homogeneity or independence?
          Circle your answer and explain why based on the study design.

          \ansline{Test for homogeneity. The researcher drew three separate random samples (one per school) and measured one categorical variable (extracurricular activity) in each group.}

    \item \textbf{Write the hypotheses.} State $H_0$ and $H_a$ in the context of this study.

          $H_0$: \ansline{The distribution of extracurricular participation (Sports/Arts/Neither) is the same across Schools A, B, and C.}
          $H_a$: \ansline{The distribution of extracurricular participation is not the same across all three schools (at least one school differs).}

    \item \textbf{Independence condition.} How was independence of observations achieved
          in this study? Was the 10\% condition necessary? Explain.

          \ansline{A stratified random sample was drawn (100 students randomly selected from each school). The 10\% condition applies if school populations are less than 1000 students each.}

    \item \textbf{Large-counts condition.} Calculate the expected counts for two cells:
          School A/Sports and School C/Neither. Then check whether the large-counts
          condition is plausibly met.

          Expected count (School A / Sports): \ans{$(100 \times 135) \div 300 = 45$}

          Expected count (School C / Neither): \ans{$(100 \times 65) \div 300 \approx 21.67$}

          Large-counts condition satisfied? \ans{Yes} \quad Explain: \ansline{All expected counts will be at least $(100 \times 65)/300 \approx 21.67 > 5$; condition is met.}

          \begin{teachernote}
              All expected counts for Scenario A: each row total is 100.
              Sports: $100 \times 135/300 = 45$ (all three schools the same).
              Arts: $100 \times 100/300 \approx 33.33$ (all three).
              Neither: $100 \times 65/300 \approx 21.67$ (all three).
              Smallest is 21.67, well above 5.
          \end{teachernote}

\end{enumerate}

\vspace{0.08in}

\begin{scenariobox}[Scenario B: Political Affiliation and Policy Views]{sky}
A random sample of 300 adults is asked their political affiliation (Democrat, Republican,
or Independent) and their view on a new policy (Support or Oppose).
The results are shown below.

\vspace{0.08in}
\renewcommand{\arraystretch}{1.4}
\begin{tabularx}{0.72\linewidth}{l ccc}
    & \textbf{Support} & \textbf{Oppose} & \textbf{Row Total} \\
    \hline
    \textbf{Democrat}    & 85 & 15  & 100 \\
    \textbf{Republican}  & 30 & 70  & 100 \\
    \textbf{Independent} & 55 & 45  & 100 \\
    \hline
    \textbf{Col Total}   & 170 & 130 & 300 \\
\end{tabularx}
\end{scenariobox}

\begin{enumerate}[label=\textbf{B\arabic*.}]

    \item \textbf{Identify the test type.} Is this a test for homogeneity or independence?
          Circle your answer and explain why based on the study design.

          \ansline{Test for independence. One random sample of 300 adults was drawn, and two categorical variables (political affiliation and policy view) were measured on each individual.}

    \item \textbf{Write the hypotheses.} State $H_0$ and $H_a$ in the context of this study.

          $H_0$: \ansline{There is no association between political affiliation and policy views in the adult population.}
          $H_a$: \ansline{There is an association between political affiliation and policy views in the adult population.}

    \item \textbf{Independence condition.} How was independence of observations achieved
          in this study? Was the 10\% condition necessary? Explain.

          \ansline{Data came from a random sample of adults (SRS). The 10\% condition applies: $300 \le 10\%$ of all adults --- easily satisfied since there are hundreds of millions of adults.}

    \item \textbf{Large-counts condition.} Calculate the expected counts for two cells:
          Democrat/Support and Republican/Oppose. Then check whether the large-counts
          condition is plausibly met.

          Expected count (Democrat / Support): \ans{$(100 \times 170) \div 300 \approx 56.67$}

          Expected count (Republican / Oppose): \ans{$(100 \times 130) \div 300 \approx 43.33$}

          Large-counts condition satisfied? \ans{Yes} \quad Explain: \ansline{All expected counts are at least $(100 \times 130)/300 \approx 43.33 > 5$; condition is met.}

          \begin{teachernote}
              Scenario B expected counts: each row total is 100.
              Support column: $100 \times 170/300 \approx 56.67$ (all three rows).
              Oppose column: $100 \times 130/300 \approx 43.33$ (all three rows).
              Smallest is 43.33, well above 5.
          \end{teachernote}

\end{enumerate}

\vspace{0.08in}

\begin{extensionbox}
\textbf{Tier E --- Extension:} Both Scenario A and Scenario B involve two categorical
variables appearing in a two-way table. A classmate says both scenarios call for a test
for independence. Is the classmate correct? Explain in 2--3 sentences what feature of
the study design determines which test to use.

\ansline{The classmate is incorrect for Scenario A. The key is how data were collected: Scenario A used three separate random samples (one per school), so we compare distributions across groups --- a test for homogeneity. Scenario B used one random sample with two variables measured on each person --- a test for independence. The presence of two variables in a table does not automatically indicate a test for independence.}
\end{extensionbox}

\end{document}
